% Authorship: Claude-drafted, author-revised
\section{The Problem of Practice}

Everything to this point has been apparatus: a model, its parameters, its
requirements, and a long series of measurements taken against them. It comes
to this. The relocation reported in Part VI is a claim about the kind of
indeterminacy a practice-based theory of legal content can carry, the
difficulty it names is not Berman's alone, and it leaves a defender a choice
among three substantive positions, none of them yet taken by principled
positivism.

\subsection{Discharging the Condition}

Part VI stated the relocation conditionally and named the two conditions.
If the divergence among the mechanisms holds across the family of
aggregation rules and across the determinacy threshold, then the
relocation is not a feature of one rule or one threshold but of the theory.
On the points that were his to win, both conditions have now been answered,
and the answers ran in Berman's favor.

As Part VII established, no rule inside the family Berman endorses is a
candidate for dissolving the divergence. In that family every rule computes
a principle's contribution as its weight multiplied by its activation, so a
disagreement about weights reaches the verdict by construction. Of the
whole battery, the only rules that could blunt it are the ones that discard
the weights outright, and discarding them abandons the weighing the theory
makes central. As Part VIII established, the divergence was never bought with a gate. The
headline is computed on the bare sign of net force with no threshold applied,
and raising the threshold increases the rate at which the mechanisms disagree
rather than reducing it, so the reported figure is the most conservative the
sweep produces.\footnote{\textit{Supra}
Part~VII.D (every rule inside the cardinal family responds to magnitude,
at rates from $0.023$ to $0.075$, against a ceiling of $0.0024$ outside
it, and every such rule computes each principle's contribution as weight
times activation); \textit{supra} Part~VIII.D (Part~VI's headline divergence is computed on the
sign of net force with no threshold applied, as is every pairwise rate that
Part reports; the four-way rate rises with the threshold, from $36.1\%$ at
$\tau = 0$ through $44.5\%$ at the default to $47.8\%$ at $\tau = 0.3$ on
correlated activations).}

With the condition discharged, the relocation is a claim about the theory
rather than about one rule or one threshold. From here on the argument draws on results already reported rather than
adding to them, with one exception noted where it arises. For a theory of
legal content, those results mean four things. The relocated indeterminacy has a shape of its own, and the
shape is not Hart's. The difficulty it names belongs to any theory that
asks practice for quantities rather than for criteria, so it is not
Berman's alone. A defender is left with three positions and no neutral
ground among them. And the whole of it invites objections, the first of
which comes from inside the theory.

\subsection{More Determinate and Less Settled}

On Hart's account the indeterminacy has a shape, and the shape matters as
much as the fact of it. It is \textit{local}: the law runs out in identifiable hard
cases, at the penumbra of a rule whose core is settled, and elsewhere it
does not run out at all. It is also \textit{visible}. Where the law has
run out, the officials applying it can tell that it has.\footnote{Hart,
\textit{supra} note 1, at 124--36 (open texture and the inevitability of
discretion in hard cases); \textit{id.}\ at 272 (Postscript).} Neither
feature is incidental. Because Hart's gaps can be located, Dworkin could object to them, and there
is no objecting to an indeterminacy nobody can point at.\footnote{Dworkin, \textsc{Law's Empire}, \textit{supra}
note 4, at 37--43.} Inside a Hartian gap, a judge exercising discretion knows she is exercising
it, and so the discretion thesis is a claim about judicial candor as much as
about legal content. For Dworkin's complaint and Hart's defense alike, the gap has to be
locatable.

On both counts the indeterminacy measured here has the opposite shape. It
is \textit{global} in a particular sense, and the sense needs care. The
divergence is not spread evenly across the law. It concentrates on the cases that turn on the moral principles, which is
where the construals are furthest apart, and the objections below return to what follows from that.

It lacks a boundary a deciding judge could draw from inside her own
construal. Under uptake the moral principles are multiplied by zero
before any judgment about them is consulted, so a case that turns on
dignity does not present to her as a case that turns on dignity. In her construal, the marker that would identify the affected set has
already been discarded. A judge is not wholly without a warning, but the warning is unreliable. Under
every construal the size of a case's net force predicts divergence better than
chance, so the closest cases are also the likeliest to come out differently
under a rival reading.\footnote{This is the one measurement first reported here rather than earlier, and it
re-reads Part VI's divergence battery rather than adding a run. Taking the
size of net force as a predictor of four-way divergence, the area under the
curve is $0.636$ under uptake, $0.711$ under reliance, $0.757$ under
appointment, and $0.901$ under strategic, against $0.5$ for a coin. In the
lowest decile of net force the four structures disagree on between $0.53$ and
$0.89$ of cases; in the highest, on between $0.000$ and $0.158$.} But the warning is construal-relative, and sharply asymmetric. Among the cases
each construal ranks in the top decile of its own net force, a reliance judge
runs a reversal risk of $0.006$ and an uptake judge $0.158$, a factor of
twenty-eight. Of the four, the construal that zeroes the moral principles is arguably the
one whose confidence is worth least, and nothing inside it says so.

The first half of that is good news for Berman. At the Article's default
gate of $\tau = 0.2$, each of the three mechanisms carried through the
threshold sweep returns a determinate verdict in a large majority of
cases: uptake on $86.4\%$, strategic on $87.7\%$, reliance on $88.5\%$,
against the stipulated engine's $86.2\%$. On every mechanism the margin over the hierarchical Hart is positive, and the
margins sit near $+0.27$: uptake $0.271$, strategic $0.283$, reliance $0.291$,
and appointment $0.284$. Two conditions travel with that. The margins are
measured on independent activations, and Part IV reports that the comparison
reverses under the correlated activations used for this Part's divergence
figures, the tiered rule there deciding more by consulting less.
And the four figures are closer together than four independent measurements
would be, since the hierarchical Hart never reads the weights, so each margin
is that mechanism's own determinacy rate less one shared constant.\footnote{Thirty seeds at $\tau = 0.2$, $800$ cases
per seed; determinacy rates $0.862$ (stipulated), $0.864$ (uptake),
$0.877$ (strategic), $0.885$ (reliance). Margins over the hierarchical
Hart at the same setting: uptake $0.271$, strategic $0.283$, reliance
$0.291$, appointment $0.284$, with the stipulated engine between $0.268$
and $0.272$ across batteries (\textit{supra} Part~IV.B, Part~VI.A,
Part~VII.B). The appointment mechanism's weights come from a separate
simulation that the threshold sweep does not rebuild, so no determinacy
rate is reported for it here; its margin comes from Part~VI's own
battery. \textit{See supra} Part~VIII.D.} Berman promised more law than
Hart. On every construal of practice tested, he delivers it.

Now the second half. On the battery of fifty thousand synthetic cases
reported in Part VI, the four weight structures return different verdicts
on $36.1\%$ of cases, and on $27.2\%$ they return opposite
\textit{determinate} verdicts once the default gate is
applied.\footnote{The divergence figures come from a different battery
and a different threshold setting than the determinacy rates above, and
the two are not two decompositions of one run. The determinacy rates are
thirty-seed means at $\tau = 0.2$ over $800$ cases per seed; the
divergence rates are computed over a single common battery of fifty
thousand cases under each mechanism's mean weight vector, with the
$36.1\%$ taken on the bare sign of net force ($\tau = 0$) and the
$27.2\%$ requiring each of a pair to clear $\tau = 0.2$. On independent
activations, the Article's conservative baseline, the corresponding rates
are $31.8\%$ and $19.9\%$. \textit{See supra} Part~VI.C, Part~VIII.D.}
Here the pairing matters. Taken alone, each construal of practice decides between eighty-six and
eighty-nine cases in a hundred, and reports every one of them the same way,
as a net force with a sign. The four of them part
company on more than a third. Individually, they are not indecisive. They
are decisive against one another.

For the argument that follows, the missing signal is the point. Reasoning under one construal of practice, a judge gets a number and a
verdict. She does not get a flag. On the
\textit{Obergefell}-style activation profile of Part VI, nothing in uptake's
output records that reliance, reading the same bench and the same holdings by
a different measure with its own warrant in Berman's text, would have decided
the case the other way. Uptake's number, $-1.154$, arrives with no annotation attached. Only by
running a second construal does the contingency become visible, and that is
an operation the theory gives nobody any reason to perform, because the
theory does not say there is a second construal to run. It says the weights are practice-grounded, and
``practice,'' in the singular, sounds like one thing.\footnote{The claim
is about the model's output. Whether a real judge has other signals
available, a sharp dissent, a lopsided circuit split, a visible drift in
the doctrine, is an empirical question about adjudication that this
Article does not reach. What the model shows is that the trace the divergence does leave in the
quantity the theory identifies as legal content, the size of net force, is
construal-relative and worth much less under some construals than others.}

So the finding, at its strongest: a theory can be more determinate and
less settled at once. On every mechanism, by a margin near $+0.27$, Berman genuinely delivers more
law than Hart, and nothing here suggests he delivers it wrongly. Because the design contains no standard against which any of them could be
wrong, no verdict reported in this Article has been shown to be the wrong
verdict. It has been shown, rather, that the determinacy is \textit{indexed}. Under a
construal it is determinacy, and the theory does not fix the construal whose
determinacy it is.

That claim needs one distinction it has not yet been given. It is not the
claim that there is no fact about which dynamic a given legal system runs on.
There may well be such a fact, and a defender who says so is not contradicted
by anything here. The claim is about the theory: given a complete description
of a practice, principled positivism does not by itself say which feature of
that practice fixes the weights, and so does not by itself determine legal
content. That is a claim about the theory's completeness rather than about the
metaphysics of law, and it is the weaker of the two claims available. Since a
simulation of practice cannot discover which dynamic a real system runs on,
the weaker claim is also the one the evidence supports. Hart's gaps announce themselves.
This one does not.

That is what the relocation does to the debate Berman entered. In the
Hart-Dworkin exchange the currency was the \textit{quantity} of law.
Hart conceded that his theory ran out in hard cases, Dworkin said that
running out in hard cases is running out where it counts, and every
subsequent move was an argument about how much law a source-based account
can deliver. Berman answers that charge, and the margin over the hierarchical Hart is
the answer. What the four mechanisms surface is not a quantity of law at
all but a property of the law delivered, and it is arguably a property
neither party had a name for, because neither theory could produce it. 

On one side a Hartian gap is the theory declining to answer, and it declines
out loud. On the other a Dworkinian right answer is the theory answering, in a way
that invites the objection that the answer is the interpreter's. A verdict that is determinate under one reading of a theory's own grounds and
reversed under another reading the theory equally licenses is a third thing,
and from the inside it looks like the second.

Two things keep that from being a larger claim than the evidence supports.
First, about two-thirds of cases are decided the same way by all four
structures, so the claim is not that the choice of construal decides the
law generally; on the run of cases practice speaks with one voice however
it is read, and what is at issue is a large and substantively serious
minority. 

Second, the margin over the hierarchical Hart is real, holds on every
mechanism tested, and is not taken back by anything said here. Part IV's finding
stands as reported there. The guard that closed Part VI holds as
well. Nothing here impugns the synchronic grounding relation: officials'
taking-up of a principle fixes its weight. Under every mechanism that relation
is intelligible, and under every mechanism it works. The open question is the
dynamics of the changing grounds, and it is likely the openness in the
dynamics that propagates into the verdicts.

\subsection{Criteria and Magnitudes}

The difficulty is not Berman's alone. To say why means saying what his theory asks of practice that Hart's does
not.

For its part, Hart's rule of recognition asks officials to converge on
\textit{criteria} of validity: enacted by this legislature, decided by this court, ratified
by this procedure, contained in this text. The convergence it asks for is qualitative. Downstream of a shared
criterion, officials can disagree about nearly everything, including which
way it cuts in the case in front of them, and the rule of recognition is
satisfied so long as they agree on what counts as a source.\footnote{Hart, \textit{supra}
note 1, at 94--99, 100--10. Hart's own account is not wholly free of
quantity: how much convergence among officials is enough to sustain a
rule of recognition is a question he leaves at the level of ``general''
acceptance, and the Introduction to this Article flags it as one of the
places his framework is silent. The contrast drawn here is accordingly
one of degree in how much quantification a theory demands of practice,
not a bright line between a qualitative theory and a quantitative one.}

Berman asks practice for something else. His weights are magnitudes:
quantities of legal force on a scale, one for each principle, and no
aggregation runs without them. A weight is not a verdict on whether a principle belongs to the system. It
is a number that says how much the principle counts when it collides with
another. 

That is a different demand on the same social fact, and it is strictly harder.
The entailment runs one way. In fixing a magnitude a practice thereby fixes the criterion, since a
principle carrying any weight above zero is a principle that counts and
nothing further need be asked. A practice can fix the criterion and leave the
magnitude wide open, since officials can agree without residue that
dignity is among the considerations a court may properly invoke while
disagreeing, or saying nothing at all, about how much it weighs against
plain text. Magnitude is criterion plus a number. That is the sense in
which the second demand is harder, and not any empirical claim that
officials converge less readily on numbers than on categories; the
empirical claim would need evidence this Article does not supply.

On magnitudes, then, is where the four mechanisms disagree. All four run on the same nine principles and the same category map. None of them proposes a
principle the others have not heard of, and none strikes one from the list.
They differ on how much each principle counts. One
qualification. In the limit a quantity of zero is a disagreement about
membership, and under uptake the four moral principles carry weight
exactly zero. But those zeros are post-cut. The underlying uptake shares are small rather than absent, running from
about $0.35$ to about $0.50$ for the moral
principles against $0.79$ to $1.00$ for the textual and structural ones. At the level of the practice the mechanisms actually
read, the disagreement is about how much, and the disagreement about
membership is downstream of a line the model draws.\footnote{\textit{Supra}
Part~VI.C (the uptake zeros are the weights after the cut separating live
principles from dead ones; on the raw shares the moral tier still carries roughly half the weight of the
textual one, and the worked case
still comes out the other way); \textit{supra} Part~V.A (where the cut
falls governs how many principles count as live).}

Stephen Perry's objection to aggregative accounts was a different one. His
was that the messy and contested practice of diverse judges cannot deliver
weights precise enough for aggregation to yield determinate results.\footnote{Perry, \textit{supra} note 9, at 794; Berman,
\textit{Practices}, \textit{supra} note 7, at 328--30 (the ``rough
weights suffice'' reply).} Berman's reply was that rough weights suffice, and Part IV vindicated him:
the weighted sum absorbs weight error on the order of a tenth without
losing its properties. 

To the difficulty here, though, that reply does not reach. It is not that any
single reading of practice delivers weights too imprecise to use. Each of the four delivers weights precise enough to clear noise tolerance, the
one requirement that tests precision. Part VI reports that it is also the only
one of the four requirements that discriminates among weight vectors at all. It is that
the readings deliver \textit{different} weights, and the distance between
them is not imprecision of the kind the model was shown to tolerate.
Against a tolerance for weight error near $0.10$, measured on the same
$[0,1]$ scale as the weights themselves, uptake assigns dignity a weight of
zero where reliance assigns it $0.95$. That is not error around a value.
It is a different value.

With the generalization come its limits. On this reading any
practice-based positivism that moves from criteria to quantities would
inherit this problem, and the more finely a theory quantifies, the more of
it it inherits. Ask practice only whether a source counts, and you need it determinate on a
yes or a no, and the several intelligible ways of reading one practice will
often agree on that. Ask practice how much a
consideration weighs and you need it determinate on a real number, which
the several intelligible readings will not deliver even where they agree
on the yes or the no. That is arguably what the four mechanisms display: agreement on the roster,
disagreement on the column of figures beside it.

Three limits attach.

First, the claim reaches only theories that quantify. It is not a refutation
of positivism, and it is not an argument that positivism ought not to
quantify, since Part IV showed that quantifying buys the extra law: the margin over the hierarchical Hart is the
payoff for weighing a plurality of principles against a rule that decides
at its highest applicable tier, and a theory that declined to quantify
would forfeit it. The margin should not be attributed to magnitude alone. Against a consensus
Hart the weighted sum wins at every bar above a bare majority, by $+0.270$,
$+0.589$ and $+0.786$ as the bar rises, and draws even only at bare-majority
counting, where the rival decides the most cases. So quantifying does buy an
advantage over counting, everywhere but at the counting rule's most permissive
setting; it does not buy the advantage over a tiered rule by
magnitude alone, since the tiered rule differs from the weighted sum in
plurality as well.\footnote{\textit{Supra} Part~IV.B (the
consensus-Hart curve runs $-0.021$ at a fifty-percent bar, then $+0.270$,
$+0.589$ and $+0.786$ as the bar rises); \textit{supra} Part~VII.A (the
hierarchical Hart differs from the weighted sum in magnitude and in
plurality at once, so no comparison against it localises the margin to
either alone). What magnitude buys is measured separately in Part~VII.C,
as a rate of disagreement with counting rules rather than as a
determinacy payoff.}

Second, Hart is not wholly exempt. His account needs an answer to how much
convergence is enough, which is a quantity, so the difference is one of degree
in how much a theory asks practice to fix.

Third, the argument establishes a demand rather than a failure. Showing that a
theory needs more from practice than its predecessor did is not showing that
practice cannot supply it. It shows where the next argument has to be made.

\subsection{What a Defender Owes}

Since the aggregation works, a defender of principled positivism owes not a
repair to it but an addition. In the first place the theory
supplies a synchronic relation: officials take a principle up and embed it in their
reasoning and their holdings, and their doing so confers its weight. It does
not supply a \textit{diachronic} component, an account of what drives those
takings-up over time. Named that way, the missing piece
sits where the evidence puts it and nowhere else. The four mechanisms are four candidate diachronic accounts, and the
divergence among them is the cost of the theory's silence about which is
right.

One reply comes before any account of the dynamics, and it is the strongest
thing a defender can say here. If the grounding relation is synchronic, a
defender may hold that a principle's weight is fixed by current takings-up,
and that the historical dynamic that deposited the present practice is a
question in legal history with no bearing on legal content. That reply
concedes everything Part VI measured and denies that a diachronic account is
owed at all.

It relocates the question rather than dissolving it. The four mechanisms are
not four histories converging on one present practice. They are four readings
of what a taking-up is, and they disagree about the weights of a single bench
and a single body of holdings taken as they stand. Uptake reads how often a
principle is invoked; reliance reads what still rests on it. A defender holding that
current takings-up fix weight has still to name the current facts that count
as takings-up, and that is the question again in its synchronic form.

Granting that an account is owed, three positions are available. Either the
dynamics do not come apart where it matters, or they do; and if they do,
either one of them governs legal content or none does. Those branches give the
three below, taken in the order a defender is likeliest to reach for them.
Each is a substantive commitment and each has a price. I name them and do not
choose.

\textit{Single out one dynamic}, and defend it as the one real legal
systems run on. Being the closest of the four to Berman's own account, reliance is the
natural candidate, and a defender who took this route would be saying that
weight accrues to principles that prevail in contested disputes and embed
themselves in holdings later relied upon, and not otherwise. The cost is
that it converts a conceptual claim into an empirical bet. At least the bet is a payable one, since the dynamics would leave different
signatures: the mean weight vectors stand at rank
correlations of $-0.78$ and $-0.75$ between uptake and each of the
holdings-reading mechanisms and $+0.93$ within that pair, and real
doctrine either exhibits that pattern or does not.\footnote{\textit{Supra}
Part~VI.B (the rank correlations); \textit{supra} Part~III.A (reliance is
the closest of the four to Berman's own account). Part~X sets out what
testing the bet would require.}

\textit{Deny that the dynamics come apart where it matters}, by arguing
that they converge in actual practice in ways they do not in simulation:
real benches, real dockets, and real doctrines of precedent may bring
invocation and embedding into line even though nothing in the models
forces them together. The cost is that the convergence has to be shown rather than assumed, and
Part VI supplies both halves of the relevant evidence. On one side there is a real convergence, reliance and appointment disagreeing
on $1.8\%$ of verdicts, though the two are not independent constructions and
Part VI reports that their agreement tracks the culture they share. On the
other there is a real opposition, uptake against reliance at $35.3\%$ and
uptake against appointment at $33.6\%$. The reply is available, and it is not free. It also carries a discount that
Part VI itself supplies: the shared moral lean behind the convergence is partly a
category-size effect, since the moral tier holds four principles to the
structural tier's three and the textual tier's two, which limits how much
weight a general convergence thesis can rest on that
result.\footnote{\textit{Supra} Part~VI.C (pairwise disagreement rates on
correlated activations: uptake-reliance $35.3\%$, uptake-appointment
$33.6\%$, uptake-strategic $20.6\%$, strategic-reliance $16.3\%$,
strategic-appointment $14.7\%$, reliance-appointment $1.8\%$);
\textit{supra} Part~VI.B (the category-size concession, and the reading
that the convergence tracks a shared object, holdings, rather than a
shared construction).}

\textit{Accept that they come apart}, and with it the indexing: weights
are indexed to a practice, and legal content is correspondingly indexed.
The cost is a form of relativity about legal content that positivists have
generally wanted to avoid, and there is a specific tension with the answer
to Dworkin. Berman's determinacy margin was offered as the reply to the charge that
positivism delivers too little law. A determinate
answer that holds under one construal of practice and reverses under
another is not obviously the determinate answer Dworkin was asking
for.\footnote{Indexing legal content to a practice is already positivism
and no concession at all; what is new is indexing it to a \textit{construal}
of a single practice, which is the distinction Section IX.E takes up
first. The concession, if made, would attach to the dynamics of the
changing grounds and not to the synchronic grounding relation, which
nothing in this Article disturbs. \textit{See supra} Part~VI.D.} It may
nonetheless be the honest position.

I decline to choose among the three, and the refusal is not evasion. Each
is a substantive commitment about how legal systems in fact work, and a commitment of that kind is the theory's to make. A model can price the options. It cannot pick one.
Picking one would require knowing the dynamic that actually drives the
takings-up in a real legal system, and that is not a question a simulation
of practice can answer about the practice it simulates.

\subsection{Objections}

The first objection is the hardest, and it comes from inside the theory. It
says that what has been reported is just positivism being sensitive to social
facts. A Bermanian will put it this way: of course legal content depends on
which practices obtain, and that is positivism. Where officials take up dignity the law differs from where they do not, and
the fact that four different practices yield four different bodies of law is
the theory working, not the theory failing. Reporting it as a difficulty
confuses the social contingency of legal content, always asserted by
positivism, with an indeterminacy in the theory, something else entirely.

The objection has force only if the four mechanisms model four social
worlds. Whether they do is a question about what the divergence is measured
on, and it needs a careful answer, because the four are not built alike.

Two concessions come first. The mechanisms are separate simulations rather
than four read-outs of one simulated history, and their internal machinery
differs more than an earlier draft of this Part allowed. Uptake decides in
nine-judge panels drawn from four methodologies. Reliance has no bench at
all, resolving its disputes against a fixed exogenous pressure structure.
Appointment replaces cohorts of judges who carry no methodology, and
strategic learns a shared weight vector. Nor do the four draw their case
streams alike: uptake's stream is independent across principles, while the
other three are category-correlated.\footnote{\textit{See supra} Part~III.A
for the four constructions. In the implementation, the uptake bench draws
its stream from a uniform distribution over each principle independently,
while reliance, strategic and appointment draw category-correlated
activations.} These four are not one bench observed four ways, and nothing
below assumes they are.

Because the argument does not require the mechanisms to share their
construction, the concession does not reach it. The argument requires the
divergence to be measured on shared cases, and it is. Each mechanism yields
a weight vector. The divergence rates of Part VI are obtained by taking
those four vectors and running each over one common battery of fifty
thousand cases. At the point of measurement the cases are identical and
only the weights differ. The four mechanisms are instruments of different
construction, and the divergence reports their readings of the same object.

That distinction is the one the objection collapses. A Bermanian who says that
different practices yield different law is describing two legal systems whose
officials behave differently, and about that there is no dispute. Part VI
reports something else: four readings of one set of cases, disagreeing about
what the law in those cases is. Contingency of legal content across worlds is
ordinary positivism and embarrasses no one. Underdetermination of legal
content within a world is what a theory of legal content would need to rule
out, and it is what the divergence measures.

Across the four, Berman's account of the ground of weight does not vary.
His schema stands in all of them: officials take principles up and
embed them, and that confers weight. The mechanisms vary in the takings-up
they count and in how those accumulate over time. That is a question about the
dynamics rather than about the grounding relation.
Uptake tallies invocation events over a history. Reliance accrues on
victories and decays. Appointment reads what survives the turning over of
the bench. Strategic settles wherever a divided bench converges. Four ways
of running one relation forward, and they do not end in the same place.

Three things support that reading. The first is the design. In Part II the weights
were held fixed while the activations varied across two Justices, which is
Berman's own experiment: Burger and Powell read the same statute
differently, and the model locates their disagreement in the activations.
In Part VI the mirror image ran, holding one activation profile fixed and
varying only the construal of practice. Same case, opposite experiment.
Run through the four, the \textit{Obergefell}-style profile returns
$-1.154$ under uptake, $+0.708$ under reliance, $-0.140$ under strategic, and $+0.509$ under
appointment. Here the claim is literal rather than figurative. Those four numbers are one activation profile multiplied by four weight
vectors. The case did not change, the principles did not change, and the activations were held fixed by construction. Only the reading of
``taking-up'' moved. Whatever one thinks of the four simulations that produced the weight
vectors, the worked case is a single object read four ways, and the readings
disagree.

The second is what governs the convergence. Reliance and appointment agree
closely, at a rank correlation of $+0.93$ and a verdict disagreement of
$1.8\%$, and they are not independent instruments: on a given seed both take
their exogenous pressure structure from the same source, and the appointment
mechanism was built with the reliance mechanism in view. What matters is not
that they agree but what their agreement tracks. Part VI decouples the
cultures, so that each reads its own rather than a common one, and the
correlation falls from $+0.83$ to $+0.30$, its dispersion across seeds rising
fourfold.\footnote{\textit{Supra} Part~VI.B. The decoupling test is a separate
run, and its baseline is the $+0.83$ mean of eight per-seed correlations
rather than the $+0.93$ computed between the mean weight vectors.} Hold the
culture fixed and the two readings converge; give them different cultures and
they come apart. That is the pattern a world-and-construal reading predicts
and a four-worlds reading does not, since on the latter there is no separate
culture to hold fixed.

The third is that the gap can be named. If Berman's text picked a construal,
the objection would also be sound. It does not. It says principles are
constituted by the taking-up behaviors of the system's subscribers, and
that the weights are practice-grounded, and it points at both invocation
and embedding without resolving in favor of either; Part VI put the open question in that form. On this reading the Bermanian reply presupposes an answer the theory has
not given.\footnote{What the four mechanisms share is narrower than an earlier draft
claimed: the nine-principle set, the category map, the aggregation rule,
and the common battery of cases on which the divergence rates are computed.
They do not share a bench, a panel structure, or the activation
distribution of the streams from which their weights are formed. The shared
battery is what the argument uses. \textit{Supra} Parts~III.A, VI.A. The worked case is
the \textit{Obergefell}-style profile of Part~VI.C, returning $-1.154$
(uptake), $+0.708$ (reliance), $-0.140$ (strategic), $+0.509$
(appointment); perturbing every activation in it with Gaussian noise,
uptake still opposes both reliance and appointment on $100.0\%$ of draws
at $\sigma = 0.05$, $99.1\%$ at $\sigma = 0.1$, and $87.5\%$ at $\sigma =
0.2$, twenty thousand draws at each level. One concession is
owed: the four mechanisms are separate simulations, not four read-outs
of a single simulated history, so the defensible general claim is that
they construe one \textit{kind} of practice under a shared institutional
setup. It is the worked case, where a single fixed activation profile is
run through four weight vectors, that makes ``one world'' literal, and
the reliance-appointment convergence that licenses reading the four as
construals rather than as worlds.}

A second objection accepts all of that and reaches for the obvious remedy. Why
not simply use reliance? A defender might take reliance, the mechanism Part III calls the closest of the four to
Berman's own account, and let it settle the weights. Two things follow, and
the second probably matters more.

The first is arithmetic. Drop uptake from the comparison and the divergence shrinks without going
away: the largest disagreement between
any two of the remaining structures is $16.3\%$, where the largest with
uptake in the set was $35.3\%$. The relocation survives with a smaller
coefficient. For a theory of what determines legal content, one case in six is still
arguably a serious share of a caseload, and Part VI reported the same
figure as the discount on its own headline.

The second is that choosing reliance \textit{is} the first of the three options set out above. It is a substantive commitment about which dynamic real legal systems run
on, and the theory has not made it and would have to defend it. That is no criticism of the choice, which may
well be the right one. It is a description of what the choice costs.
Picking reliance answers the problem of practice. It does not avoid it.

A third objection says the Hartian validation rule is a strawman, and its
force should be conceded first. A Hartian will say that the rule of
recognition identifies \textit{sources} of law, not conflicting substantive
standards, and that a rule so constructed would never certify both sides of
a single case in the first place. Asking only whether each authority is
valid and which way it points, the validation rule of Part IV is not a rule
any careful Hartian would propose, and a result about its behavior is not a
result about Hart.

Three answers, in ascending order of importance. In framing the conflict problem for positivism, Shapiro and Greenberg set
up that foil, and it is used here as they use it, to state a difficulty a
source-based account has to answer rather than to describe a position
anyone holds. It is offered throughout as a contradiction foil and never
as exegesis of Hart, and Part IV says so where it is introduced. And no determinacy figure anywhere is taken against it. Every determinacy comparison runs against the charitable
\textit{hierarchical} Hart, the rule sorting principles into tiers and
deciding a case at its highest applicable tier, with a consensus Hart reported alongside as a
curve across agreement bars. The margin that Part IV claims, and that the argument
above relies on, is the margin over the hierarchical Hart, and it would
stand unchanged if the validation rule were struck from the Article
entirely.

The concession that follows from all three is one of scope rather than of
substance. The contradiction result of Part IV is a result about a rule of a certain
construction, not one about Hart's theory, and this Article does not claim
otherwise. Its claim is narrower and survives the
objection: a source-based account owes some answer to the case in which
two standards both satisfy its criteria and point opposite ways, and the
weighted sum supplies one that a rule of that construction cannot. A
Hartian who says the rule of recognition would never get into that
position is describing a different answer to the same question, not
showing that the question is idle.\footnote{Shapiro, \textsc{Legality}, \textit{supra} note 6, at
270--75; Greenberg, \textit{How Facts Make Law}, \textit{supra} note~42, at
185--89.
Against the hierarchical Hart the stipulated engine's margin is $+0.272$,
positive on every seed; against the consensus Hart the margin runs from a
dead heat at bare-majority counting to $+0.270$ at a sixty percent bar,
$+0.589$ at seventy, and $+0.786$ at eighty. \textit{Supra} Part~IV.B.}

A fourth objection is that the principle set is stipulated, and it is live and
partly unanswered. As implemented, the nine
principles and the sizes of their tiers, two textual, three structural,
four moral, are the author's choice and not a finding. Part VI already concedes that the
shared moral lean under reliance and appointment is partly a category-size
effect: the moral tier is the largest, so under the correlated
activations of realistic cases it accrues broader aggregate force simply
by being larger. The global sensitivity analysis of Part VIII varies four
inputs, the threshold, the dispersion of the weights, the correlation
among activations, and the share of principles carrying no weight, and
the principle set is not among them. And the divergence runs straight
through the moral tier: uptake zeros it, reliance and appointment lift it
highest, strategic parks it in the middle, so the cases most exposed to
the choice of construal are the morally freighted ones.

Those facts interact in a way that ought to be uncomfortable. The moral tier is the largest of the three, and it is also the tier on
which the mechanisms stand furthest apart. So the same design choice that lets that tier accrue broader aggregate
force under correlated activations may also govern how much divergence
there is to measure. Because the sensitivity analysis of Part VIII does not
vary the principle set, it cannot separate the two effects. A replication that
varies the tier sizes directly can bound their joint contribution.

Restricting the principle set directly is the natural test, and it bounds the
effect. Take the four mean weight vectors and drop columns until the tiers are
equal, two principles apiece, and the four-way divergence falls from $0.361$
to $0.303$. Leave the structural tier at three and cut only the moral tier to
two, and it falls to $0.249$. So tier composition accounts for part of the
divergence, between a sixth and a third of the four-way rate, and not for the
rest: on every restricted set the structures still disagree on roughly a
quarter to a third of cases.

Two limits on that test. It restricts rather than replicates. The weight
vectors were themselves formed under the nine-principle map, so dropping
columns bounds what tier size contributes to the verdicts and leaves untouched
what it contributed to the weights. And a balanced set of three principles a
tier does not exist here, since the textual tier holds only two, so the
balanced variants are six and seven principles rather than nine and are not
like-for-like with the full set.\footnote{Tier composition as implemented:
textual (text, original meaning), structural (precedent, structure,
federalism), moral (dignity, equality, liberty, welfare). The restricted sets
are computed on correlated activations and on the bare sign of net force, as
the $0.361$ headline is; no gated counterpart was run. A second balanced
variant returns $0.297$. The sensitivity inputs are listed \textit{supra}
Part~VIII.B, and the principle set is not among them. The category-size
concession is \textit{supra} Part~VI.B, and the concentration of the
divergence in the moral tier \textit{supra} Part~VI.C.}

A fifth objection is that four mechanisms are only a sample, and it is
conceded. Four construals are a sample of the readings a practice-based theory
could take of its own key term rather than the space of them, and nothing
here rules out a fifth that behaves differently from all four. The limit
cuts both ways. A larger sample could add structures to the three already found. Since each of the three is exhibited by a mechanism that has been built and
run, it could not remove them. The sampling limit bounds
the completeness of the claim, not its truth.

The last objection is that the cases are synthetic, which is conceded and was
flagged at the outset. The cases are synthetic and the practice is simulated. A model of this kind can show what a
theory's structure supports and leaves open; the construal of
practice that describes real adjudication is an empirical question, and no number in this Article answers it, nor would any number drawn from
synthetic cases.

\subsection*{What Remains}

The problem of practice is real, and it is narrow. It is not that
principled positivism fails. The aggregation is sound, the weights can emerge from practice rather than
waiting on a theorist, and the margin over the hierarchical Hart, measured on
independent activations, survives every mechanism tested. The problem is that a theory asking practice for magnitudes has to name the
reading of practice that supplies them. Until it does, its determinacy is
determinacy-under-a-construal, and its verdicts carry a contingency that
nothing in them reveals. That is a smaller and more
tractable problem than the one principled positivism inherited from Hart
and Dworkin: not a question about whether social facts can ground legal
content, but a question about which facts about a single practice do the
grounding, with a measurable consequence attached to each answer.
